% Source: https://artofproblemsolving.com/community/c3h1962184_base_10_and_17
% Stable identity: aops-topic-1962184-q01
Write $s_{10}(n)$ and $s_{17}(n)$ for the two digit sums. Since $10\equiv1\pmod9$ and $17\equiv1\pmod{16}$,
\[
n\equiv s_{10}(n)\pmod 9,\qquad n\equiv s_{17}(n)\pmod{16}.
\]
If $s_{10}(n)=s_{17}(n)=S$, then $n-S$ is divisible by $9$ and by $16$, hence by $144$. For $n>10$ we have $S<n$ (a number with at least two digits exceeds its digit sum), so $n-S\ge144$, i.e. $n\ge144+s_{10}(n)$.

Check $145\le n\le 153$, which requires $s_{10}(n)=n-144$:
\begin{itemize}
\item $n=145,\dots,149$: $s_{10}=10,\dots,14$, never equal to $n-144=1,\dots,5$.
\item $n=150$: $s_{10}=6=n-144$, but $150=8\cdot17+14$, so $s_{17}=8+14=22\ne6$.
\item $n=151,152$: $s_{10}=7,8=n-144$, but $151=8\cdot17+15$ and $152=8\cdot17+16$ give $s_{17}=23,24$.
\item $n=153=9\cdot17+0$: $s_{17}=9$ and $s_{10}=1+5+3=9$.
\end{itemize}
The answer is $\boxed{153}$ (base $17$: $90_{17}$). Checked by computer.
